\subsection{A central-lift obstruction for the \texorpdfstring{$2D2$}{2D2} rules}
\label{subsec:no-central-lift-2d2}
\begin{proposition}[No central invertible lift for the $2D2$ rules]
\label{prop:no-central-lift-2d2}
Let $\mathcal Q$ be a pseudo-unitary fusion category with simple objects
$1,f,r,s$ and fusion rules
\[
 f^2=1,\qquad fr=rf=s,\qquad
 r^2=s^2=1+2r+2s,\qquad rs=sr=f+2r+2s.
\]
There is no invertible object $z\in Z(\mathcal Q)$ whose underlying
object is $f$.
\end{proposition}
\begin{proof}
Put $d=2+\sqrt5$, so $d^2=1+4d$. The displayed rules imply
$\mathcal Q_{\rm ad}=\mathcal Q$. Consequently the only central
lift of the tensor unit is the tensor unit: such lifts are tensor
automorphisms of the identity, equivalently characters of the universal
grading group. Thus a hypothetical $z$ satisfies $z^2\cong1$.
In its pointed braided subcategory the positive ribbon twist
$q=\theta_z$ belongs to $\{1,-1,i,-i\}$.

If $q=1$, then $1\oplus z$ is the regular algebra of a central
Tannakian $\operatorname{Rep}(C_2)$. Its module category has two simple
free orbits, represented by $1,r$, and their rule is
$\bar r^2=1+4\bar r$. This contradicts the rank-two fusion-category
classification~\cite{Ostrik2}.

Suppose $q=\pm i$. Write $I$ for induction to the center.
The formal-codegree induction formula, as recalled in
\cite[Appendix~A.2]{EIP}, gives the distinct constituents of $I(1)$
dimensions
\[
 1,\quad d^2,\quad1+2d,\quad1+2d.
\]
Every constituent has twist $1$. Each has tensor-unit multiplicity one
after forgetting, and none contains $f$, because $F I(1)$ has no
$f$-summand. The projection formula gives
$I(f)\cong z\otimes I(1)$, disjoint from $I(1)$.
Since $z^2=1$, the monodromy of $z$ with any simple central object is
$\pm1$. Hence every constituent of $I(f)$ has twist $\pm i$.

In $I(r)\oplus I(s)$ the constituents belonging to $I(1)$ have total
dimension
\[
 K=4d^2+4(1+2d)=8+24d.
\]
Indeed their total multiplicities are, respectively, $0,4,2,2$.
The constituents belonging to $I(f)$ have the same total dimension.
Since
\[
 \dim\bigl(I(r)\oplus I(s)\bigr)
 =2d(4+8d)=16+72d,
\]
the other constituents have total dimension $24d$. The first-indicator
identity gives zero twist trace on $I(r)\oplus I(s)$. Its $I(1)$ part
has real trace $8+24d$, its $I(f)$ part has purely imaginary trace,
and the absolute real trace of the remaining part is at most $24d$.
This is impossible.

Finally suppose $q=-1$. Let $\mathcal I$ be an Ising category with
fermion $\psi$ and noninvertible $\sigma$. The object
$z\boxtimes\psi$ is a central boson in
$\mathcal Q\boxtimes\mathcal I$. Condensing its regular algebra gives
a pseudo-unitary fusion category $\mathcal D$ of rank six. All twelve
simples before condensation have free orbits. Representatives give
objects $1,f,u,r,s,v$ of dimensions
\[
 1,1,\sqrt2,d,d,\sqrt2d,
\]
with the above rules on $1,f,r,s$ and additional rules
\[
 \begin{gathered}
 fu=u,\quad fv=v,\quad u^2=1+f,\quad ur=us=v,\\
 uv=r+s,\quad rv=sv=u+4v,\quad v^2=1+f+4r+4s.
 \end{gathered}
\]
We prove directly that these rules have no pseudo-unitary
categorification.

Put $D=\dim(\mathcal D)=8+16d$ and $T=Dd=16+72d$.
The six formal codegrees show that $I(1)$ has two constituents of each
dimension $1,d^2,2+4d$, all of twist $1$, each with tensor-unit
multiplicity one. Positivity and linear independence over
$\mathbb Q(\sqrt5)$ show that their restrictions have no $u$ or $v$
summands. The two constituents of dimension $d^2$ restrict to
\[
 1+a_jr+(4-a_j)s\qquad (0\le a_j\le4),
\]
and the two of dimension $2+4d$ restrict to
\[
 1+f+b_jr+(4-b_j)s\qquad (0\le b_j\le4).
\]
The $r$-coefficient of $FI(1)$ is $8$, so if $a=a_1+a_2$ then
$b_1+b_2=8-a$. The total dimension $K$ of the twist-one summands of
$I(r)$ supplied by $I(1)$ is therefore
\[
 K=a d^2+(8-a)(2+4d)=16-a+32d\ge8+32d.
\]
Let $R=T-K$ be the dimension of the other summands. The first two
indicator identities~\cite{NgIndicators}, with $\nu_2(r)\le1$, imply
\[
 \sum_j w_j\cos t_j=-K,\qquad
 \sum_j w_j\cos(2t_j)=D\nu_2(r)-K\le D-K,
 \qquad \sum_j w_j=R,
\]
where $w_j>0$ and $e^{it_j}$ are their dimensions and twists.
If $R=0$ the first equation is already impossible; otherwise
Cauchy--Schwarz gives
\[
 2K^2\le R(R+D-K),\qquad
 0\le T(T+D)-K(3T+D).
\]
But the last expression is at most
\[
 T(T+D)-(8+32d)(3T+D)
 =-64-512d-1088d^2<0.
\]
This final contradiction excludes $q=-1$.
\end{proof}

\begin{remark}
The published Galois pseudo-unitarity reduction~\cite[Lemma~3.1]{SchNear}
for categories having
two Frobenius--Perron dimensions upgrades this statement to arbitrary
complex categorifications of the displayed ring: a central lift survives
Galois conjugation. This uses the corrected published formulation
``Galois conjugate to a pseudo-unitary category'', not a claim that every
categorification is itself pseudo-unitary.
\end{remark}

\begin{remark}[Scope of the inner-conjugation construction]
Let $\mathcal P=\mathcal Q_{\rm pt}$ and
$\mathcal D=\mathcal Z_{\mathcal P}(\mathcal Q)$.
By~\cite[Remark~2.4(i)]{GNN}, this relative center is Morita equivalent
to $\mathcal P\boxtimes\mathcal Q^{\rm rev}$.
Its invertibles are precisely those of $\mathcal Z(\mathcal P)$,
whose group is $C_2\times C_2$ for either associator on
$\mathcal P\simeq\operatorname{Vec}_{C_2}^{\omega}$.
Thus this natural inner-conjugation construction cannot realize the
$C_4$-based rings $R8\_1\_9$ or $R8\_1\_18$; no indicator assumption
is needed.

The displayed rank-four rules have trivial universal grading, so the
theorem implies $\operatorname{Inv}\mathcal Z(\mathcal Q)=1$.
Morita invariance then gives
$|\operatorname{Inv}\mathcal Z(\mathcal D)|=4$.
If $\mathcal D$ had the rules $R8\_1\_49$, its universal grading
would also be trivial. The forgetful map on invertibles would therefore
be injective, since its kernel is the group of tensor automorphisms of
the identity, and hence bijective. Every invertible would have a central
lift, contradicting the noncommuting invertible action in that ring.
This rules out the proposed construction only. General exclusions are
proved independently in Theorems~\ref{thm:r9-center-trace}
and~\ref{thm:r18-sts-moments} and
Proposition~\ref{prop:r49-fs-morita}.
\end{remark}
